
% =====================================================================
\section{Cosmology I: can the model resolve the Hubble tension?}
\label{sec:hubble}
% =====================================================================

Hypothesis H3 has two parts. This section examines the first, the Hubble tension, and
Sec.~\ref{sec:dm} the second, the dark matter. Together they meet objective O6. Everything was
computed with a program written for this project, an Einstein-Boltzmann solver that follows the
expansion of the universe, recombination and the linear perturbations of all species; the physics of
the model enters through two additions to it, a dark fluid of radiation with a step (this section)
and dark matter that is a superfluid (Sec.~\ref{sec:dm}).

\subsection{What the model contains, cosmologically}
\label{sec:darksector}

The four-dimensional model of Sec.~\ref{sec:vs} has, besides the Standard Model, a \emph{dark
sector}: the fermions $\psi$, with 16 degrees of freedom, eight internal states each of particles and
antiparticles, and the mass $m$; the radial mode, with one degree of freedom and the mass $m_s$; and
the massless phase mode, with one. They interact with one another with the strength $\gVS$, and with
ordinary matter only through gravity; any further coupling would be an addition to the model. A dark
sector that was populated in the early universe has a temperature of its own, $T_d$. Its history
depends on the fermion mass, and two regimes must be distinguished.

\emph{Light fermions}, with $m$ between about 1 and 100 eV, are relativistic until late. As long as
$T_d\gg m$ the sector is a tightly coupled gas of radiation: the Yukawa interaction keeps it in
equilibrium, and it does not stream freely. When $T_d$ falls below $m$ the fermions and antifermions
annihilate into phase quanta. Entropy is conserved, so that the phase quanta are heated, and the
energy density of the dark radiation rises relative to that of radiation that simply redshifts. This
\emph{step} is the mechanism of the ``stepped dark radiation'' that Aloni \emph{et al.} proposed for
the Hubble tension \cite{Aloni2022}. If the radial mode is heavier than the fermions, the gas has one
bosonic and sixteen fermionic degrees of freedom before the step and one bosonic after it, and the
step factor is
\begin{equation}
r_{\rm step}=\Big[\frac{1+\tfrac78\cdot16}{1}\Big]^{1/3}=15^{1/3}=2.466,
\label{eq:rstep}
\end{equation}
much larger than the $(15/7)^{1/3}=1.289$ of the Wess-Zumino model of those authors.

\emph{Heavy fermions}, with $m$ above about a keV, annihilate long before the epochs that the
microwave background probes. What remains is the phase quanta, a small amount of radiation, and, if
there is an excess of fermions over antifermions, the excess fermions, which are non-relativistic
matter and pair up as described in Secs.~\ref{sec:cjt} and \ref{sec:order}: a candidate for dark
matter.

A single field has a single mass, so that \emph{the model cannot be in both regimes at once}. With
fermions of an eV it has interesting dark radiation and no dark matter; with heavy fermions it has a
candidate for dark matter and almost no dark radiation. The two parts of hypothesis H3 are therefore
tested separately.

\subsection{The Einstein-Boltzmann solver and its validation}
\label{sec:solver}

The program follows the standard treatment of linear cosmological perturbations in the synchronous
gauge \cite{MaBertschinger1995}: photons with polarization, massless neutrinos, baryons, cold dark
matter, a cosmological constant and the dark fluid. Recombination is computed with the rate equations
of the effective three-level atoms of hydrogen and helium
\cite{SeagerSasselovScott1999,SeagerSasselovScott2000,WongMossScott2008}, the helium abundance from
standard nucleosynthesis including the effect of extra radiation \cite{Planck2015XIII}, and the spectra by integration
along the line of sight \cite{SeljakZaldarriaga1996}: the temperature (TT), polarization (EE) and
cross (TE) spectra of the microwave background to multipole 2500, and the matter power spectrum.
Appendix~\ref{app:num} gives the details. The \emph{fiducial model} throughout is the standard
\LCDM{} model with the densities $\omega_b=0.02237$ and $\omega_c=0.1200$ of baryons and of cold dark
matter, the Hubble constant $H_0=67.36~\kmsMpc$, the amplitude $A_s=2.1\times10^{-9}$ and tilt
$n_s=0.9649$ of the primordial spectrum, the optical depth $\tau=0.0544$, and massless neutrinos with
the standard effective number of species, 3.044.

The solver was validated against spectra of the fiducial model computed with the established code
\textsc{class} \cite{BlasLesgourguesTram2011}. Those spectra include gravitational lensing, which
the solver does not compute; to compare like with like, its spectra were lensed with the lensing
potential of the reference by the correlation-function method \cite{LewisChallinor2006}.
Figure~\ref{fig:cmbval} shows the differences, and Table~\ref{tab:derived} compares derived
quantities. For multipoles between 100 and 2000 the two codes agree within {{GEN:cmb_max_tt}} per
cent in TT, within {{GEN:cmb_max_te}} per cent of $\sqrt{D^{TT}D^{EE}}$ in TE, and within
{{GEN:cmb_max_ee}} per cent in EE; the matter power spectrum agrees within {{GEN:pk_max}} per cent
for wavenumbers between $10^{-4}$ and 3 Mpc$^{-1}$. Halving all tolerances and all steps in
wavenumber changes the spectra by less than {{GEN:cmb_conv_max}} for multipoles above 100. This
agreement was not obtained at once: the comparison exposed five errors of the first version of the
program, which are recorded, with the tests that now guard against them, in the error log of the
project \cite{SpacetimesRepo}.

\begin{figure*}
\includegraphics[width=\textwidth]{fig_cmb_validation}
\caption{\label{fig:cmbval}Differences between the solver of this project and \textsc{class} for the
fiducial model. Dashed: spectra without lensing; the oscillating difference is the smoothing of the
acoustic peaks by lensing, which \textsc{class} includes. Solid: spectra lensed with the lensing
potential of the reference.}
\end{figure*}

\begin{table}
\caption{\label{tab:derived}Derived quantities of the fiducial model: the angular size $\theta_s$ of
the sound horizon, the sound horizon $r_{\rm drag}$ at the end of the baryon drag epoch, the amplitude
$\sigma_8$ of the matter fluctuations, and the redshifts of last scattering, of the drag epoch and of
reionization.}
\begin{ruledtabular}
\begin{tabular}{lcc}
quantity & this work & \textsc{class} \\
\hline
{{ROWS:cmb_derived_table}}
\end{tabular}
\end{ruledtabular}
\end{table}

\subsection{The mechanism, and the measure of success}

The microwave background measures the angle $\theta_s=r_s/D_A$ under which the sound horizon $r_s$ at
recombination is seen at the distance $D_A$. In the \LCDM{} model, $\theta_s$ and the densities fix
$H_0=67.4~\kmsMpc$ \cite{Planck2018VI}; the local distance network gives $73.50\pm0.81$
\cite{H0DN2026}. Extra radiation before recombination speeds up the expansion and shrinks $r_s$; to
keep $\theta_s$, $D_A$ must shrink too, which means a larger $H_0$. The difficulty is that extra
radiation also changes the damping of the small-scale anisotropies and the phases of the acoustic
peaks, which the data constrain. The measure of success used here is therefore: \emph{how much does
the fit to the microwave background worsen when the model is forced to give a certain $H_0$?}

\subsection{The dark radiation of the model}

The dark fluid is described by its density today, $N$, in units of the density of one massless
neutrino species, and by the transition redshift $z_t$, defined by $m=T_{d0}(1+z_t)$, where $T_{d0}$
is the temperature of the dark radiation today. Its equation of state and sound speed follow from the
thermodynamics of the mixture of one massless boson and sixteen fermionic degrees of freedom of mass
$m$ at constant entropy (Fig.~\ref{fig:step}). The fluid is tightly coupled: it has perturbations of
density and velocity but no anisotropic stress, in contrast to free-streaming neutrinos.

\begin{figure*}
\includegraphics[width=\textwidth]{fig_step_thermo}
\caption{\label{fig:step}The stepped dark fluid of the model (solid) and, for comparison, the
Wess-Zumino fluid of Ref.~\cite{Aloni2022} (dashed). (a) Energy density relative to radiation that
only redshifts, normalized to late times: before the step the density is lower by the factor
$1/r_{\rm step}$. (b), (c) During the step the pressure and the sound speed drop below the value
$1/3$ of pure radiation.}
\end{figure*}

Seven one-parameter families of models are compared; in each, the density $N$ is varied. Family 0 is
free-streaming radiation, additional massless neutrinos. Family 1 is a fluid without step, the limit
$z_t\to\infty$: the phase quanta of a model with heavy fermions, if they interact enough to be a
fluid. Families 2, 3 and 4 are the fluid of the model with its step in the epochs that the microwave
background probes, $z_t=10^4$, $2\times10^4$ and $5\times10^4$. Family 5 is the Wess-Zumino fluid
with $r_{\rm step}=(15/7)^{1/3}$ and $z_t=2\times10^4$, the best-fit region of Ref.~\cite{Aloni2022}.
Family 6 is the fluid of the model with an early step, $z_t=10^6$: after nucleosynthesis, and before
the scales of the microwave background enter the horizon. For $N=0.4$ the transition redshifts
$10^4$, $10^{5.5}$ and $10^6$ correspond to the fermion masses {{GEN:dr_mass_zt_4}},
{{GEN:dr_mass_zt_5p5}} and {{GEN:dr_mass_zt_6}} eV.

\subsection{Mock data and method}

The ``observed'' spectra are those of the fiducial \LCDM{} model, without a noise realization, so
that \LCDM{} itself fits with $\chi^2=0$ and every $\chi^2$ below is a difference $\Delta\chi^2$
relative to \LCDM. Two data sets are simulated. The first, \emph{Planck-like}, covers 60 per cent of
the sky with the noise and the beams of the three channels of the Planck satellite at 100, 143 and
217 GHz, multipoles from 30 to 2500 in temperature and to 2000 in polarization, and a prior on the
optical depth in place of the large-scale polarization. The second adds a \emph{deep survey}: 25 per
cent of the sky observed with a beam of 1.4 arcminutes and a noise of 10 $\mu$K arcmin, the
sensitivity of the sixth data release of the Atacama Cosmology Telescope. The likelihood is the exact
likelihood of Gaussian fields on a fraction of the sky. For each family and each $N$ the six
parameters of the \LCDM{} model are refitted; a fit alternates exact computations of the spectra with
Newton minimizations of the exact likelihood of spectra linearized with stored derivatives, and the
reported $\chi^2$ is always that of an exact computation. A test with derivatives recomputed at the
end point of a fit bounds the error of this procedure by {{GEN:dr_chord_max}} in $\chi^2$ and by
{{GEN:dr_chord_max_h0}} $\kmsMpc$ in $H_0$.

\begin{table}
\caption{\label{tab:fisher}One-standard-deviation errors of the six parameters of the \LCDM{} model
from the Fisher matrix of the two sets of mock data, the Planck-like one (P) and the one with the
deep survey (P+D), compared with the published errors of the Planck collaboration
\cite{Planck2018VI} (for $h$: the error of $H_0$ divided by 100).}
\begin{ruledtabular}
\begin{tabular}{lccc}
parameter & mock P & Planck 2018 & mock P+D \\
\hline
{{ROWS:lcdm_fisher_table}}
\end{tabular}
\end{ruledtabular}
\end{table}

Does the mock resemble the real experiment? Table~\ref{tab:fisher} compares the errors of the
\LCDM{} parameters predicted for the Planck-like mock with those published by the Planck
collaboration; they agree within 25 per cent for every parameter. For free-streaming radiation the
mock gives an error of {{GEN:sigmaN_free_planck}} on $N$, where the Planck collaboration found
$N_{\rm eff}=2.92\pm0.19$ \cite{Planck2018VI}: the mock is more constraining by a quarter, mainly
because it contains no lensing. All values of $\Delta\chi^2$ below should be read with this in mind;
for the real Planck data they would be smaller by a factor of about {{GEN:dr_mock_factor}}, the
square of the ratio of the two errors.

\subsection{Results}

\begin{figure*}
\includegraphics[width=\textwidth]{fig_dr_profiles_planck}
\caption{\label{fig:dr}Planck-like mock data, six of the seven families (the family with
$z_t=5\times10^4$ is omitted for clarity; it is in Table~\ref{tab:dr}). (a) Worsening of the best
fit, relative to \LCDM, as a function of the density $N$ of the extra radiation. (b) The Hubble
constant of the best fit; the band is the local measurement. (c) The same information combined: the
cost of a given Hubble constant.}
\end{figure*}

\begin{table*}
\caption{\label{tab:dr}Results for the Planck-like mock data. $\sigma(N)$ and $dH_0/dN$ (in
$\kmsMpc$ per unit of $N$) are from the Fisher matrix at the fiducial point; the other columns are
from the exact fits: the density $N$ and the Hubble constant at $\Delta\chi^2=4$, the Hubble constant
at $\Delta\chi^2=9$, the density $N_{73.5}$ at which the best-fit $H_0$ equals the local value 73.5,
and its cost. An asterisk marks a linear extrapolation beyond the largest fitted $N$.}
\begin{ruledtabular}
\begin{tabular}{lcccccccc}
family & $\sigma(N)$ & $dH_0/dN$ & $N$ at $\Delta\chi^2=4$ & $H_0$ there & $H_0$ at $\Delta\chi^2=9$ & $N_{73.5}$ & its $\Delta\chi^2$ & $\sqrt{\Delta\chi^2}$ \\
\hline
{{ROWS:dr_summary_table_c0}}
\end{tabular}
\end{ruledtabular}
\end{table*}

Figure~\ref{fig:dr} and Table~\ref{tab:dr} give the results for the Planck-like mock data.

\emph{Every family can reach the local $H_0$, and every one at a prohibitive price.} The cheapest is
the fluid of the model with an early step, at $\sqrt{\Delta\chi^2}=$ {{GEN:dr_sig_local_c0_f6}}; the
fluid without step and the Wess-Zumino fluid need {{GEN:dr_sig_local_c0_f1}} and
{{GEN:dr_sig_local_c0_f5}}, and free-streaming radiation {{GEN:dr_sig_local_c0_f0}}.

\emph{A step of the model that falls into the epochs probed is the most expensive case.} Its large
step leaves a strong imprint: modes that enter the horizon before the step see less radiation than
modes that enter after it, which changes the spectrum as a function of the multipole in a way that
the other parameters cannot absorb. For $z_t$ between $10^4$ and $5\times10^4$ the local value costs
$\sqrt{\Delta\chi^2}=$ {{GEN:dr_sig_local_min_model}} to {{GEN:dr_sig_local_max_model}}.

\emph{An early step is slightly cheaper than no step, because of helium.} If the step is complete
before the scales of the microwave background enter the horizon, these scales see a fluid of constant
density $N$; but during nucleosynthesis the fluid had only the density $N/r_{\rm step}$, and less
helium was made. At $N=0.5$ the helium mass fraction is {{GEN:dr_yhe_early}} instead of
{{GEN:dr_yhe_nostep}}. Less helium means more free electrons and less damping of the small-scale
anisotropies, which offsets part of the additional damping caused by the radiation; $\Delta\chi^2$
falls from {{GEN:dr_chi2_half_nostep}} to {{GEN:dr_chi2_half_early}}.

\emph{At a cost of $\Delta\chi^2=4$ the gain is small:} $H_0=$ {{GEN:dr_h0_2sigma_c0_f6}} for the
early step, {{GEN:dr_h0_2sigma_c0_f1}} for the fluid without step, {{GEN:dr_h0_2sigma_c0_f0}} for
free-streaming radiation, and between {{GEN:dr_h0_2sigma_min_model}} and
{{GEN:dr_h0_2sigma_max_model}} for a step in the epochs probed. At best this is
{{GEN:dr_gain_2sigma_best}} $\kmsMpc$ of the 6.1 that are needed.

\begin{figure*}
\includegraphics[width=\textwidth]{fig_dr_zt}
\caption{\label{fig:zt}The fluid of the model with $N=0.4$ after the step: cost and Hubble constant
as functions of the transition redshift. A step long before the epochs probed by the microwave
background acts like a fluid without step that was less dense during nucleosynthesis; a step after
recombination means that only $N/r_{\rm step}$ was present before recombination, so that both the
cost and the gain are small; in between, the step is seen by the data and is expensive.}
\end{figure*}

Figure~\ref{fig:zt} shows the dependence on the transition redshift for $N=0.4$. An early step,
$z_t\gtrsim10^{5.5}$, costs $\Delta\chi^2=$ {{GEN:dr_zt_early_chi2}} and gives $H_0=$
{{GEN:dr_zt_early_h0}}; a step after recombination, $z_t\lesssim10^3$, costs
{{GEN:dr_zt_late_chi2}} and gives {{GEN:dr_zt_late_h0}}; the most expensive case, at
$z_t=10^{{{GENM:dr_zt_worst_g}}}$, costs {{GEN:dr_zt_worst_chi2}}. With the deep survey everything
tightens: the local value of $H_0$ then costs $\sqrt{\Delta\chi^2}=$ {{GEN:dr_sig_local_c1_f6}} for
the early step, {{GEN:dr_sig_local_c1_f1}} for the fluid without step and
{{GEN:dr_sig_local_c1_f0}} for free-streaming radiation.

\subsection{The phase mode as a fractional neutrino, and published analyses of real data}

Even with heavy fermions the model has one unavoidable piece of dark radiation, the massless phase
mode. If the dark sector was once in thermal contact with ordinary matter and decoupled when the
plasma had $g_*$ effective degrees of freedom, a massless boson contributes today
\cite{Weinberg2013}
\begin{equation}
\Delta N_{\rm eff}=\frac47\Big(\frac{43}{4\,g_*}\Big)^{4/3},
\label{eq:weinberg}
\end{equation}
which is {{GEN:dr_weinberg_lo}} for a decoupling above the electroweak scale, $g_*=106.75$, and
{{GEN:dr_weinberg_hi}} for a decoupling just before the muons annihilate, $g_*=57/4$. With the slope
$dH_0/dN$ of free-streaming radiation in Table~\ref{tab:dr} such an amount raises $H_0$ by at most
{{GEN:dr_weinberg_gain}} $\kmsMpc$: a welcome but insufficient shift, and the largest value is itself
disfavored by the data quoted next.

The simulation assumes that the sky is like the \LCDM{} model and asks what the model then costs.
Whether the real data prefer extra radiation is an empirical question that others have answered. The
Planck collaboration found $N_{\rm eff}=2.92\pm0.19$ for free-streaming radiation, with no preference
for more \cite{Planck2018VI}. The Atacama Cosmology Telescope, combined with Planck, found
$N_{\rm eff}=2.86\pm0.13$ and, for self-interacting dark radiation, the family 1 of this paper,
$N_{\rm idr}<0.134$ at 95 per cent confidence \cite{Calabrese2025}. With the slope of
Table~\ref{tab:dr} that limit allows a gain of {{GEN:dr_act_gain}} $\kmsMpc$. For the stepped fluid of
Ref.~\cite{Aloni2022} and a second interacting model, the data of that telescope restrict the density
about as strongly as for simple models \cite{CvetkoJosephMarquesTavares2025}. The limits loosen to
$\Delta N_{\rm eff}<0.58$ if the primordial spectrum is allowed a running tilt and a running of the
running \cite{GarnyNiedermannSloth2026}, a caveat that applies equally to the fixed power-law
spectrum used here. Our mock results agree with this picture: for the fluid without step the forecast
limits at $\Delta\chi^2=4$ are $N<$ {{GEN:dr_n_2sigma_c0_f1}} from the Planck-like mock and $N<$
{{GEN:dr_n_2sigma_c1_f1}} from the mock with the deep survey, to be compared with the published
limit of 0.134 from the real data of both experiments.

\subsection{Verdict on the Hubble tension}

The dark radiation of the model raises $H_0$, and its stepped form is exactly the kind of fluid that
was proposed for this purpose. But, first, for data consistent with the \LCDM{} model the density
needed for the local value of $H_0$ is excluded at $\sqrt{\Delta\chi^2}\ge$
{{GEN:dr_sig_local_min_all}} in every family simulated, and at $\sqrt{\Delta\chi^2}\ge$
{{GEN:dr_sig_local_min_model}} if the step of the model falls into the epochs probed. Second, the
published analyses of the real data find no preference for such radiation, and limits near
$N\approx0.13$. Third, the step of the model, $15^{1/3}$, is larger than that of the Wess-Zumino
model and is \emph{more} visible to the data, not less, unless it is complete before the scales of
the microwave background enter the horizon; and then the model is the fluid of family 6, which still
fails by {{GEN:dr_sig_local_c0_f6}} standard deviations. \emph{Hypothesis H3 is not supported for the
Hubble tension: with the amount of fluid dark radiation that the published limit allows, the model
removes about {{GEN:dr_act_fraction}} per cent of the discrepancy.} This agrees with the conclusion
of the earlier study for bound neutrino fields \cite{NashNCSS}, now reached for a different dark
sector with an independent code.

\subsection{What has been established, and what has not}

Established: the solver agrees with \textsc{class} to the precision stated, and its results are
converged. For data of the quality of the Planck satellite that follow the \LCDM{} model, the cost of
raising $H_0$ with each of seven kinds of dark radiation, including the stepped fluid that the model
predicts, and its dependence on the transition redshift.

Not established, or assumed: the mock analysis is not an analysis of the real data. It has no
lensing, no foregrounds, an idealized likelihood and a power-law primordial spectrum, and its limits
on $N$ are tighter than those of Planck by about a quarter; the conclusions about the real sky rest
on the published analyses cited, with which the mock results are consistent. The step was simulated
for a radial mode much heavier than the fermions; if it is lighter there are two smaller steps
instead of one. The dark radiation was assumed to be a perfect fluid at all times, whereas after the
step the phase quanta interact only through derivative couplings and may begin to stream freely;
families 0 and 1 bracket this. Finally, the fermion mass was taken to be constant through the step.
By Sec.~\ref{sec:thermal} the mass melts at a temperature $T_{1/2}$ between
{{GEN:cjt_thermal_thalf_m_lo}} and {{GEN:cjt_thermal_thalf_m_hi}} fermion masses, and the assumption
holds only if $T_{1/2}\gg m/3$, which requires $\gVS\lesssim0.3$; for a stronger coupling the step
would coincide with that transition, which was not simulated.
